% Provenance: integral 2249, chapter 24, pp. 609--618; capitulo_24_dualidad_t.tex.
% Operator-domain formulas are made explicit on the realization already constructed.
\section{T-duality: integer coordinates, change of section, and operator equivalence}
\label{iv:dualidad}

The arithmetic decomposition of APP preserves the residual representative
and its quotient. TPK evolution also transports its orientation and
the memory of boundary crossings. Interchanging the two
coordinates requires preserving this distinction and specifying the domain
on which the operation is again defined. An arithmetic involution is
first constructed on that domain; its quadratic form
subsequently admits a representation by self-adjoint operators.

\subsection{Stable domain of the interchange}

In the positive residual section,
\begin{equation}
 r_9^+(n)=1+((n-1)\bmod9),\qquad
 Q_9^+(n)=\frac{n-r_9^+(n)}9,
 \qquad n=r_9^+(n)+9Q_9^+(n).
 \label{iv:eq:division-positiva}
\end{equation}
An interchange can place an arbitrary quotient in the first
coordinate. Its stable domain is therefore
\(\mathscr D_0=\ZZ^2\times\RR_{>0}\), not merely the section
\(1\leq r\leq9\). Extending the domain makes it possible to define
a total operation; the admissibility of a pair as the readout of
a TPK route separately retains its selection rules.

\begin{definition}[Normalized compact form and interchange]
\label{iv:def:dualidad}
For \(d=(m,w,R_0)\in\mathscr D_0\), let
\begin{equation}
 E_{R_0}(m,w)=\frac{m^2}{R_0^2}+w^2R_0^2,
 \qquad
 \mathsf T_0(m,w,R_0)=(w,m,R_0^{-1}).
 \label{iv:eq:energia-t}
\end{equation}
The symbols \(m,w\) name the two transported integer
coordinates; \(R_0\) is the positive dimensionless scale of
the form. Their interpretation through momentum, winding, and
physical radius is established by a subsequent realization map.
\end{definition}

\begin{theorem}[Involution and preservation of the quadratic form]
\label{iv:thm:dualidad-escalar}
The map \(\mathsf T_0\) is an involution of
\(\mathscr D_0\), and
\[
 E_{R_0^{-1}}(w,m)=E_{R_0}(m,w).
\]
\end{theorem}
\begin{proof}
Two interchanges restore \((m,w)\), and two inversions restore
\(R_0\). Moreover,
\[
 E_{R_0^{-1}}(w,m)=w^2R_0^2+\frac{m^2}{R_0^2}
 =E_{R_0}(m,w).
\]
\end{proof}

The equivalence holds for each radius and its reciprocal. At \(R_0=1\),
the radius remains fixed and the operation acts on the same form.
A specific pair is fixed by the full involution precisely
when, in addition, \(m=w\).

\subsection{Compact Hamiltonian and self-adjoint domain}

Let \(\mathcal H_{\mathrm{lat}}=\ell^2(\ZZ^2)\), with
orthonormal basis \(\{|m,w\rangle\}\). Define
\begin{align}
 H(R_0)|m,w\rangle&=E_{R_0}(m,w)|m,w\rangle,\label{iv:eq:hamiltoniano}\\
 D(H(R_0))&=\left\{\psi:\sum_{m,w\in\ZZ}
 E_{R_0}(m,w)^2|\psi_{m,w}|^2<\infty\right\}.\label{iv:eq:dominio-h}
\end{align}
Finitely supported vectors belong to the domain and form a
core for the diagonal operator. The associated quadratic form has
domain
\[
 D(H(R_0)^{1/2})=\left\{\psi:\sum_{m,w}
 E_{R_0}(m,w)|\psi_{m,w}|^2<\infty\right\}.
\]

\begin{proposition}[Self-adjointness and compact resolvent]
\label{iv:prop:h-autoadjunto}
The operator \(H(R_0)\) is self-adjoint and nonnegative. Its
eigenvalues are \(E_{R_0}(m,w)\), counted with multiplicity,
and \((H(R_0)+I)^{-1}\) is compact.
\end{proposition}
\begin{proof}
A multiplication operator defined by a real sequence is symmetric
on \eqref{iv:eq:dominio-h}. If \(\eta\) belongs to the domain
of the adjoint, testing the adjoint identity against each vector
\(|m,w\rangle\) forces its image to have coordinates
\(E_{R_0}(m,w)\eta_{m,w}\). Membership in
\(\ell^2\) is exactly condition \eqref{iv:eq:dominio-h}; hence
the adjoint has the same domain and the same action. Nonnegativity
follows by summing the nonnegative weights.

With \(c_{R_0}=\min(R_0^{-2},R_0^2)>0\), we have
\(E_{R_0}(m,w)\geq c_{R_0}(m^2+w^2)\). Outside increasing
finite sets, the coefficients
\((1+E_{R_0}(m,w))^{-1}\) tend uniformly to zero.
Truncating the resolvent to those sets produces finite-rank
operators converging in norm, which proves its
compactness. The diagonal basis determines the stated spectrum.
\end{proof}

\begin{theorem}[Unitary conjugation, including domains]
\label{iv:thm:t-unitaria}
\label{prop:dualidad-radios-k}
The interchange
\(U_T|m,w\rangle=|w,m\rangle\) is unitary,
\(U_T^2=I\), and satisfies
\begin{equation}
 U_TD(H(R_0))=D(H(R_0^{-1})),
 \qquad U_TH(R_0)U_T^{-1}=H(R_0^{-1}).
 \label{iv:eq:conjugacion-h}
\end{equation}
It also transports the quadratic-form domains.
\end{theorem}
\begin{proof}
The interchange is an involutive permutation of the
orthonormal basis. For a sequence \(\psi\), the identity in
Theorem \ref{iv:thm:dualidad-escalar} gives
\[
 \sum_{m,w}E_{R_0^{-1}}(m,w)^2|(U_T\psi)_{m,w}|^2
 =\sum_{m,w}E_{R_0}(m,w)^2|\psi_{m,w}|^2.
\]
Equality of the domains is established. On each coordinate,
\(U_TH(R_0)U_T^{-1}\) multiplies by
\(E_{R_0}(w,m)=E_{R_0^{-1}}(m,w)\). The same sum, with
the first power of the weight, proves the result for the
quadratic forms.
\end{proof}

The scalar equality and the operator equality are two realizations
of the same involution. The second preserves the domains that
allow the first to be interpreted as spectral equivalence,
not merely as an equality on finitely supported vectors.

\subsection{Specialization to the radius of the action ellipse}

The radius \(R_\star\) determined in Section
\ref{iv:ellipse:section} fixes the realization of this collection that
corresponds to the angular representations of the state. The symbol
\(R_0\) allows the transport results to be expressed for the entire
positive collection; composition with the ellipse evaluates that collection at
\(R_0=R_\star\) and also preserves its reciprocal image.

\begin{corollary}[Elliptic collection stable under duality]
\label{iv:cor:radio-especializado}
The set
\[
 \mathscr D_{\mathrm{el}}
 =\ZZ^2\times\{R_\star,R_\star^{-1}\}
\]
is stable under \(\mathsf T_0\). Its Hamiltonians are self-adjoint,
nonnegative, and have compact resolvent. The interchange \(U_T\)
transports the domain of \(H(R_\star)\) onto that of
\(H(R_\star^{-1})\), including their form domains.
\end{corollary}
\begin{proof}
The positivity of \(R_\star\) and its unique determination are established
in Section \ref{iv:ellipse:section}. Inversion permutes
\(R_\star\) and \(R_\star^{-1}\), while interchange preserves
\(\ZZ^2\). Proposition
\ref{iv:prop:h-autoadjunto} and Theorem \ref{iv:thm:t-unitaria}
apply on these two fibres.
If \(R_\star=1\), the two fibres coincide; otherwise
the operator intertwines two distinct fibres of the same collection.
\end{proof}

The common action scale determines the area of the ellipse; its semiaxis
ratio determines \(R_\star\). This separation makes it possible to preserve
the area and the common action scale while interchanging the semiaxes and obtaining
the reciprocal radius. The index interchange on \(\ell^2(\ZZ^2)\)
and the transports of the action plane act on different spaces;
their respective properties are proved on their own domains.

\subsection{Change of representative and quotient memory}

The zero section uses \(\bar r\in\{0,\ldots,8\}\). The change
that preserves the integer in \eqref{iv:eq:division-positiva} is
\begin{equation}
 C(r,Q)=(\bar r,Q+\mathbf1_{\{r=9\}}).
 \label{iv:eq:cambio-seccion}
\end{equation}
The modification of the quotient is part of the operation.
For \(n=9\), the two pairs are \((9,0)\) and \((0,1)\).
Their uncorrected energies are \(81/R_0^2\) and \(R_0^2\).
Representing the same integer therefore does not imply equality of
those two insertions into a form that distinguishes the coordinates.

The equivalence relation is preserved by defining
\[
 L_9=\ZZ(9,-1),\qquad S(x_1,x_2)=(x_2,x_1),
 \qquad q_{R_0}(x)=\frac{x_1^2}{R_0^2}+x_2^2R_0^2.
\]
For \(x\in\ZZ^2\), let
\begin{equation}
 \mathcal E_{R_0}^{L_9}([x])
 =\min_{k\in\ZZ}q_{R_0}(x+k(9,-1)).
 \label{iv:eq:energia-cociente}
\end{equation}
This construction operates on classes. The memory distinguishing
individual representatives remains in the TPK state preceding
the quotient; the function in \eqref{iv:eq:energia-cociente}
determines the invariant evaluation of the chosen class.

\begin{theorem}[Covariant duality of sections]
\label{iv:thm:dualidad-cociente}
The energy in \eqref{iv:eq:energia-cociente} is well defined.
The transformation
\[
 \mathsf T_{\mathrm{cov}}([x]_{L_9},R_0)
 =([Sx]_{SL_9},R_0^{-1})
\]
is a bijection between
\((\ZZ^2/L_9)\times\RR_{>0}\) and
\((\ZZ^2/SL_9)\times\RR_{>0}\), with inverse given by the
second interchange, and
\[
 \mathcal E_{R_0^{-1}}^{SL_9}([Sx])
 =\mathcal E_{R_0}^{L_9}([x]).
\]
\end{theorem}
\begin{proof}
Replacing \(x\) by \(x+k_0(9,-1)\) reindexes the candidates
for the minimum. The minimum exists: the function of \(k\) is quadratic and
has leading coefficient \(81/R_0^2+R_0^2>0\). The change
\((9,Q)\mapsto(0,Q+1)\) has difference in \(L_9\), so
it represents the same class. Interchange transports
the equivalence relation to \(SL_9\), and
\[
 \min_{\ell\in L_9}q_{R_0^{-1}}(Sx+S\ell)
 =\min_{\ell\in L_9}q_{R_0}(x+\ell).
\]
Finally, \(S^2=I\) recovers the initial quotient and radius.
\end{proof}

The minimum can be evaluated exactly using two
candidates. If \(x=(m,w)\), the real minimum of the quadratic
occurs at
\[
 k_* =\frac{wR_0^4-9m}{81+R_0^4}.
\]
The integer minimum is attained at \(\lfloor k_*\rfloor\) or
\(\lceil k_*\rceil\). This formula does not alter the equivalence;
it makes its finite evaluation explicit and allows the effect
of changing section to be checked without discarding the quotient.

\subsection{Radius normalization in the string realization}

Let \(\alpha'>0\) be a scale of the string realization.
This symbol denotes the string scale and is kept distinct
from the fine-structure constant \(\alpha\). Define
\begin{equation}
 \mathsf T_{\alpha'}(m,w,R)
 =\left(w,m,\frac{\alpha'}R\right),\qquad
 \mathcal N_{\alpha'}(m,w,R)
 =\left(m,w,\frac R{\sqrt{\alpha'}}\right).
 \label{iv:eq:normalizacion-cuerda}
\end{equation}

\begin{theorem}[Exact conjugation by the scale]
\label{iv:thm:normalizacion-cuerda}
Normalization satisfies
\[
 \mathcal N_{\alpha'}\mathsf T_{\alpha'}
 =\mathsf T_0\mathcal N_{\alpha'}.
\]
If
\[
 E_{\mathrm{str}}(m,w,R)
 =\frac{m^2}{R^2}+\frac{w^2R^2}{\alpha'^2},
\]
then
\(E_{\mathrm{str}}=\alpha'^{-1}E\circ\mathcal N_{\alpha'}\).
The combinations
\[
 p_L=\frac mR+\frac{wR}{\alpha'},\qquad
 p_R=\frac mR-\frac{wR}{\alpha'}
\]
transform as \(p_L\mapsto p_L\) and \(p_R\mapsto-p_R\).
\end{theorem}
\begin{proof}
Both compositions in the first equality give
\((w,m,\sqrt{\alpha'}/R)\). On the other hand,
\[
 E_{R/\sqrt{\alpha'}}(m,w)
 =\frac{\alpha'm^2}{R^2}+\frac{w^2R^2}{\alpha'}
 =\alpha'E_{\mathrm{str}}(m,w,R).
\]
Substituting \((w,m,\alpha'/R)\) into the expressions for
\(p_L,p_R\) gives the two asserted transformations,
respectively.
\end{proof}

The conjugation fixes what it means to express the same result in
a dimensional chart. In the elliptic collection, the radii are
\[
 R=\sqrt{\alpha'}R_\star,\qquad
 R^\vee=\sqrt{\alpha'}R_\star^{-1}.
\]
Their product is \(\alpha'\). The internal determination of the
dimensionless semiaxis ratio and the length-squared scale of the
realization serve different functions: the former fixes the shape;
the latter expresses its radii in a dimensional unit.

The preceding normalization allows the same result to be expressed with
normalized or physical radius. Assigning a tension and a
unit to \(\alpha'\) belongs to the realization map and has
its own origin; it is not a datum entering the production
of APP--TRIT--TPK registers.

\subsection{Relationship with the modular transformation}

Radius inversion has a precise representation on the
imaginary axis of the upper half-plane. Let
\(\iota(R_0)=\mathrm iR_0^2\) and
\(\mathsf S(\tau)=-1/\tau\). Then
\begin{equation}
 \mathsf S\circ\iota(R_0)=\iota(R_0^{-1}).
 \label{iv:eq:semiconjugacion-modular}
\end{equation}
The equality follows from
\(-1/(\mathrm iR_0^2)=\mathrm iR_0^{-2}\). For the modular
function \(J=j-744\), whose invariance is preserved in the
Moonshine development, it follows that
\(J(\mathrm iR_0^2)=J(\mathrm iR_0^{-2})\).

The map \(\iota\) expresses a semiconjugacy of parameters.
The interchange \(U_T\) acts on the coordinates of
\(\ell^2(\ZZ^2)\); \(\mathsf S\) acts on the upper
half-plane; the dual automorphism of an orbifold acts on its
graded sectors. Writing out their domains makes it possible to
compose specific relationships between them while preserving the
information of each representation. The results assembled here
establish coordinate interchange with explicit
operator domains and proofs.
